% restatement.tex — 一、问题重述与分析（2026-08-15 修订: 背景与问题分离, 分点呈现）

\subsection{问题背景}

某生鲜商超经营六大品类共 251 种蔬菜单品。蔬菜保鲜期短，当日未售出的商品品相变差，
部分品种隔日无法再售，因此商超每天凌晨依据历史销售与需求情况决定补货方案，并以
"成本加成"方式定价，对运损和品相变差的商品打折销售。商超的销售空间有限，蔬菜供应
品种随季节变化，凌晨进货时单品与批发价尚不确定，这些经营条件共同决定了补货与定价
决策需要在信息不完整、空间受约束的条件下做出。

题目给出的数据包括：附件1 的商品与品类信息，附件2 的 2020 年 7 月至 2023 年 6 月
逐笔销售流水（含是否打折），附件3 的逐日批发价格，附件4 的各单品近期损耗率。

\subsection{问题重述}

本题共包含四个问题：

\begin{enumerate}[label=（\arabic*）]
\item 分析各蔬菜品类与单品的销售量分布规律及品类间、单品间的相互关系；
\item 分析各品类销售总量与成本加成定价的关系，并制定 2023 年 7 月 1 日至 7 日各品类的
日补货总量与定价策略，使商超收益最大；
\item 将决策细化到单品层面，要求可售单品总数控制在 27 至 33 个且单品订购量不小于最小
陈列量 2.5 千克，依据 2023 年 6 月 24 日至 30 日的销售记录给出 7 月 1 日的单品补货量
与定价策略；
\item 给出商超还应采集的数据，并说明这些数据对补货与定价决策的帮助及理由。
\end{enumerate}

\subsection{问题分析}

\subsubsection{问题一的分析}

问题一是描述性统计与相关分析问题。分布规律一侧，销量数据具备时间维度，可从月度、
周内与日内三个尺度观察品类与单品的销售结构；相互关系一侧，需要先判断变量间关系的
形态再选取相关系数：若两变量仅存在单调关联而不满足线性与正态前提，斯皮尔曼秩相关
比皮尔逊相关更合适，配合显著性检验筛选出有统计意义的强关联组合，并联系蔬菜的上市
季节与食用搭配给出解释。

\subsubsection{问题二的分析}

问题二分两个层次。第一层是关系分析：成本加成定价由批发价与加成率共同决定，品类
日销量与当日定价构成成对观测，通过多项式曲线拟合可以得到"销量-定价"关系的显式
表达式，同时获得各品类需求的定价弹性信息。第二层是优化决策：未来一周的销量与批发价
均未知，需要先给出预测。销量与批发价的历史数据都存在明显的周内规律，用历史同期
同星期均值作为预测值。目标函数除正常销售外还包含损耗品打折回收，形成"正常销售收入、
打折回收收入与进货成本相抵"的利润结构；决策变量为补货量时模型是线性规划，若把价格
离散为若干档位并引入 0-1 选价变量，则扩展为混合整数规划，可得到定价与补货的联合
优化方案。两个方案互为参照：前者贴合成本加成的经营惯例，后者度量定价优化的潜力上限。

\subsubsection{问题三的分析}

问题三是问题二的单品级细化，并增加了两项约束：可售单品数 27 至 33 个、最小陈列量
2.5 千克。与品类层面不同，单品层面需要先解决"选哪些单品"的评价问题，仅凭利润单一指标
会忽略销售稳定性，因此用利润额、销售量、销售次数、打折次数与损耗率五个指标做综合
评价，筛选候选集后再以利润最大为目标建立 0-1 规划模型。题面要求尽量满足市场对各品类
蔬菜的需求，模型中加入各品类需求满足率的下限约束，使利润目标与需求均衡同时得到保障。
展示空间有限这一题设通过空间容量约束落实，容量由历史销量水平反推估计。

\subsubsection{问题四的分析}

问题四是开放的数据建议问题。现有模型对现实做了若干简化，如假设剩余商品全部打折售出、
需求仅随星期变化、加成率长期不变。论证思路为：指出每项简化假设对应的数据缺口，说明
所需数据的采集方式，给出将新数据引入模型后的目标函数或约束的修正形式，并分析修正
效果，形成从模型痛点、数据定义、形式化修正到预期效果的完整闭环。
